% Bewertungspaket zum Gesamttest «Exponential- und Logarithmusfunktionen» — Quelle des PDFs (python3 scripts/build-lp-pdf.py).
% Ganze Punkte je Teilschritt; Werte mit python3 nachgerechnet (04.10.2026, Folgefehler-Fälle eingeschlossen).
\documentclass[11pt]{article}
\usepackage{../lp-druck}
\renewcommand{\lpname}{Exponential- und Logarithmusfunktionen}
\renewcommand{\lpbereich}{Schwerpunktfach}
\renewcommand{\lpteilgebiet}{3.4}
\renewcommand{\lpfuss}{Bewertungspaket}
\newcolumntype{P}{>{\centering\arraybackslash}p{10mm}}
\newenvironment{raster}{\par\smallskip\noindent\tabularx{\linewidth}{@{}>{\raggedright\arraybackslash}p{38mm}P X@{}}\toprule
  \small\textsc{Teilschritt} & \small\textsc{P} & \small\textsc{Musterlösung}\\\midrule}{\bottomrule\endtabularx\par}
\newcommand{\fehler}[1]{\par\smallskip{\small\textcolor{lprot}{\bfseries Typische Fehler:} #1}}
\newcommand{\E}{\,{\footnotesize\bfseries\color{lpgruen}(E)}}
\newcommand{\A}{\,{\footnotesize\bfseries\color{lpblau}(A)}}
\begin{document}
\lpkopf{Bewertungspaket zum Gesamttest}{}

\begin{lpachtung}{Erst lösen, dann öffnen}
Dieses Paket enthält die Musterlösung. Wer es vor dem Lösen liest, misst am Ende nur, wie gut das Abschreiben gelingt.
\end{lpachtung}

\section*{1 · So gehst du vor}
\begin{enumerate}[leftmargin=*,itemsep=2pt]
\item \textbf{Gesamttest lösen} — vollständig, mit Rechenweg, ohne dieses Paket und ohne Taschenrechner.
\item \textbf{Fotografieren oder scannen}, gut lesbar, eine Seite pro Bild. \textbf{Kein Name} auf den Bildern.
  Handyfotos tragen oft unsichtbar Standort, Zeit und Gerät mit (Metadaten). Lade darum lieber einen Scan oder ein
  Bildschirmfoto des Fotos hoch, oder entferne vorher die Standortangaben.
\item \textbf{Bewerten lassen.} Ob du dafür eine KI nutzt und welche, entscheidest du selbst und in eigener Verantwortung —
  je nachdem, was dir aufgrund deines Alters und deiner persönlichen Situation erlaubt ist (Altersgrenzen und
  Nutzungsbedingungen des Anbieters, allenfalls Einverständnis deiner Eltern). Lade nur deine Lösung und dieses PDF hoch,
  keine persönlichen Daten, und füge den Auftrag aus Abschnitt~2 ein. Ohne KI kannst du dich mit Abschnitt~3 selbst bewerten.
\item \textbf{Rückmeldung prüfen.} Stimmt, was die KI aus deiner Handschrift gelesen hat? Eine KI kann sich irren —
  im Zweifel gilt das Raster in Abschnitt~3.
\item \textbf{Gezielt wiederholen} nach Abschnitt~4.
\end{enumerate}

\needspace{14\baselineskip}
\section*{2 · Auftrag an die KI}
\begin{tcolorbox}[colback=lplinie!25,colframe=lplinie,boxrule=0.5pt,arc=1mm,fontupper=\small]
Du bist Korrektorin oder Korrektor für Mathematik an einer Schweizer Berufsmaturitätsschule. Im Anhang findest du (1)~das Bewertungspaket zum Gesamttest «Exponential- und Logarithmusfunktionen» mit Musterlösung und Punkteraster und (2)~Fotos meiner handschriftlichen Lösung.\par\medskip
Geh so vor:\par
1. Schreib zu jeder Aufgabe G1 bis G8 zuerst ab, was du in meiner Lösung liest — Rechenweg und Ergebnis. Was du nicht sicher lesen kannst, markierst du mit [unsicher]. Nicht raten.\par
2. Vergib die Punkte genau nach dem Raster im Bewertungspaket (Abschnitt 3), ganze Punkte je Zeile. Ziehe einen Fehler nur einmal ab: Rechne ich mit einem falschen Zwischenergebnis richtig weiter, gibt es die Folgepunkte. Ausnahme: Zeilen mit \textbf{(E)} gibt es nur für das richtige Ergebnis, nie als Folgepunkt. Die Zeilen «Typische Fehler» sagen, welche Rasterzeile bei diesem Fehler 0 Punkte gibt — sie sind kein zusätzlicher Abzug.\par
3. Andere richtige Lösungswege zählen voll. Gleichwertige Schreibweisen zählen voll: Dezimalpunkt oder Dezimalkomma, Bruch oder Dezimalzahl ($\tfrac18 = 0.125$), $y = \ldots$ oder $f(x) = \ldots$, $\left(\tfrac12\right)^{t/6}$ oder $0.5^{t/6}$ oder $2^{-t/6}$, $\lg$ oder $\log_{10}$ — ausser wo das Raster eine bestimmte Form verlangt. Drei Stufen: Zeilen mit \textbf{(A)} verlangen nur Ablesen, diese Punkte gibt es auch ohne Rechenweg. Zeilen mit \textbf{(E)} sind Ergebnispunkte: Das Ergebnis muss stimmen \emph{und} aus einem erkennbaren Weg hervorgehen, ausser die Zeile sagt ausdrücklich etwas anderes. Alle übrigen Zeilen bewerten den Weg selbst.\par
4. Begründe jeden Abzug in einem Satz und nenne die Fehlerart (zum Beispiel «Prozentsatz statt Faktor» oder «linear statt exponentiell gerechnet»). Schreib die Musterlösung nicht ab — zeig mir, wo mein Fehler liegt.\par
5. Gib zum Schluss eine Tabelle «Aufgabe | Punkte | Begründung», die Gesamtpunktzahl (von 25) und — nach der Selbsteinschätzung in Abschnitt 4 — welches Kapitel des Leitprogramms ich wiederholen soll. Keine Note.\par
6. Fehlt ein Foto oder ist eine Aufgabe unlesbar, sag es, statt sie zu bewerten.
\end{tcolorbox}

\needspace{16\baselineskip}
\section*{3 · Musterlösung und Punkteraster}
\textbf{Allgemein:} Ganze Punkte je Zeile. Ein Fehler wird nur einmal abgezogen (Folgefehler). Gleichwertige Wege und
Schreibweisen zählen voll. In der Spalte P:
\textbf{(A)}~= nur ablesen, zählt auch ohne Rechenweg · \textbf{(E)}~= Ergebnispunkt, Ergebnis richtig
\emph{und} Weg erkennbar (ausser die Zeile sagt ausdrücklich etwas anderes) · ohne Zeichen = die Zeile bewertet den Weg.
Der ganze Test ist ohne Taschenrechner gestellt.
«Typische Fehler» nennt, welche Zeile 0 Punkte gibt, und die Punkte, die bleiben.
\textbf{Teilrichtige Zeilen:} Verlangt eine Zeile mehrere Angaben und stimmt nur ein Teil, gibt die Zeile 0 Punkte —
ausser die Zeile nennt ausdrücklich eine Aufteilung.

\begin{lpaufgabe}{G1}{$y = 2^{x}$ und $y = \left(\tfrac12\right)^{x}$ skizzieren}{3 P}
\begin{raster}
Skizze & 1 & $2^x$ steigend durch $\left(-1 \mid \tfrac12\right)$, $(0 \mid 1)$, $(1 \mid 2)$, $(2 \mid 4)$, $(3 \mid 8)$; $\left(\tfrac12\right)^x$ als Spiegelbild an der $y$-Achse, fallend. Beide über der $x$-Achse, keine erreicht sie. Toleranz: eine halbe Kästchenbreite.\\
Punkt und Asymptote & 1\A & gemeinsamer Punkt $(0 \mid 1)$; Asymptote $y = 0$ (die $x$-Achse). Beides nötig.\\
Wertemenge & 1\A & $W = \mathbb{R}^+$ (gleichwertig: alle positiven Zahlen, $y > 0$)\\
\end{raster}
\fehler{Eine Kurve schneidet oder berührt die $x$-Achse: Skizzen-Zeile 0 $\to$ 2 von 3\,P.
$W = \mathbb{R}$ oder $W = \mathbb{R}_0^+$ (die $0$ gehört nicht dazu): letzte Zeile 0 $\to$ 2 von 3\,P.
Beide Kurven steigend gezeichnet: Skizzen-Zeile 0 $\to$ 2 von 3\,P.}
\end{lpaufgabe}

\begin{lpaufgabe}{G2}{Basis aus einem Punkt}{3 P}
\begin{raster}
(a) Ansatz & 1 & $a^{-3} = 8 \Rightarrow \tfrac{1}{a^3} = 8 \Rightarrow a^3 = \tfrac18$\\
(a) Basis & 1\E & $a = \tfrac12$\\
(b) Basis & 1\E & $a^{3/2} = 27 \Rightarrow a = 27^{2/3} = \left(\sqrt[3]{27}\right)^2 = 9$. Probe $9^{3/2} = 27$.\\
\end{raster}
\fehler{(a) $a = 2$ (Kehrwert übersehen): Basis-Zeile 0; ist der Ansatz richtig notiert, bleibt dessen Punkt $\to$ 2 von 3\,P.
(a) $a = -2$: Basis-Zeile 0 — die Basis ist positiv.
(b) $a = 18$ ($27 : \tfrac32$ gerechnet): Zeile (b) 0 $\to$ höchstens 2 von 3\,P.
(b) $a = 27^{3/2}$ (Exponent falsch herum): Zeile (b) 0.}
\end{lpaufgabe}

\begin{lpaufgabe}{G3}{Population $N(t)$}{4 P}
\begin{raster}
(a) Modell & 1 & $N(t) = 500 \cdot 1.2^{t}$\\
(b) $N(2)$ & 1\E & $500 \cdot 1.44 = 720$ Tiere\\
(c) Prozent & 1\E & $1.2^2 = 1.44$: Zunahme um $44\,\%$ (gleichwertig: $\tfrac{720}{500} = 1.44$)\\
(c) Begründung & 1 & Im zweiten Jahr wächst die schon vergrösserte Population um $20\,\%$: $20\,\%$ von $600$ sind $120$, nicht $100$. Prozente addieren sich nicht, die Faktoren multiplizieren sich.\\
\end{raster}
\fehler{$N(t) = 500 \cdot 0.2^{t}$ (Prozentsatz statt Faktor): Modell-Zeile 0; $N(2) = 20$ als Folgefehler zählt nicht als (E) $\to$ höchstens 2 von 4\,P.
$N(t) = 500 + 100t$ (linear): Modell-Zeile 0, $N(2) = 700$: (b) 0, «40\,\%»: (c) 0 $\to$ 0 von 4\,P.
«$40\,\%$»: Prozent-Zeile 0; die Begründungs-Zeile ebenfalls 0 $\to$ höchstens 2 von 4\,P.}
\end{lpaufgabe}

\begin{lpaufgabe}{G4}{Medikament $m(t)$}{3 P}
\begin{raster}
(a) Modell & 1 & $m(t) = 240 \cdot \left(\tfrac12\right)^{t/6}$\\
(a) $m(18)$ & 1\E & $18 : 6 = 3$ Halbwertszeiten: $240 \cdot \tfrac18 = 30$ mg\\
(b) Zeit & 1\E & $240 : 15 = 16 = 2^4$: vier Halbwertszeiten, $t = 24$ Stunden\\
\end{raster}
\fehler{$m(18) = 80$ ($240 : 3$, linear gedacht): Zeile (a) zweite 0 $\to$ 2 von 3\,P.
$m(18) = 40$ ($240 : 6$): Zeile (a) zweite 0.
(b) $t = 96$ h ($16 \cdot 6$): Zeile (b) 0 — gezählt werden die Halbierungen ($2^4 = 16$), nicht der Faktor $16$.}
\end{lpaufgabe}

\begin{lpaufgabe}{G5}{Basiswechsel}{3 P}
\begin{raster}
(a) & 1\E & $8 = 2^3$, also $8^x = 2^{3x}$\\
(b) & 1\E & $\tfrac19 = 3^{-2}$, also $\left(\tfrac19\right)^x = 3^{-2x}$\\
(c) & 1\E & $e^{\ln 5} = 5$, also $e^{(\ln 5)x} = 5^{x}$\\
\end{raster}
\fehler{(a) $2^{4x}$ ($8 : 2$): Zeile (a) 0. (b) $3^{2x}$ (Minus vergessen): Zeile (b) 0. (c) $e^{5x}$ oder $(\ln 5)^x$: Zeile (c) 0.
Ergebnisse ohne einen Schritt wie «$8 = 2^3$»: Zählt voll, wenn die Potenzschreibweise der Basis erkennbar ist; reine Ergebnisse ohne jeden Zusammenhang 0.}
\end{lpaufgabe}

\begin{lpaufgabe}{G6}{Wasserkocher $f(t) = 90 - 70\,e^{-kt}$}{3 P}
\begin{raster}
(a) Start- und Sättigungswert & 1\A & $f(0) = 90 - 70 = 20$ °C; Sättigungswert $90$ °C. Beides nötig.\\
(b) Verlauf und Asymptote & 1 & steigt (Start $20 < 90$); Asymptote $y = 90$. Beides nötig.\\
(c) Temperatur & 1\E & Rückstand $70 \to 35 \to 17.5$ nach zwei Halbierungen: $f(10) = 90 - 17.5 = 72.5$ °C\\
\end{raster}
\fehler{(a) Starttemperatur $70$ oder $-70$: Zeile (a) 0 $\to$ 2 von 3\,P.
(c) $45$ °C ($90$ halbiert, dann nochmals halbiert, oder Wert statt Rückstand halbiert): Zeile (c) 0 $\to$ 2 von 3\,P.
(c) $55$ °C ($90 - 35$, nur einmal halbiert): Zeile (c) 0.}
\end{lpaufgabe}

\begin{lpaufgabe}{G7}{$f(x) = 2^{x} - 3$ umkehren}{4 P}
\begin{raster}
(a) aufgelöst & 1 & $y = 2^x - 3 \Rightarrow 2^x = y + 3 \Rightarrow x = \log_2(y + 3)$\\
(a) getauscht & 1\E & $f^{-1}(x) = \log_2(x + 3)$\\
(b) $D$ und Asymptote & 1 & $x + 3 > 0$: $D = \,]-3;\, \infty[$; Asymptote $x = -3$. Beides nötig.\\
(c) Nullstelle & 1\E & $\log_2(x+3) = 0 \Rightarrow x + 3 = 1 \Rightarrow x_0 = -2$\\
\end{raster}
\fehler{$f^{-1}(x) = \log_2(x) + 3$ oder $\log_2(x - 3)$: Zeile (a) zweite 0; (b) und (c) als Folgefehler bewerten, (c) nur, wenn der Weg zur eigenen Funktion passt, aber ohne (E)-Punkt $\to$ höchstens 2 von 4\,P.
(b) $D = \mathbb{R}^+$ (Verschiebung übersehen): Zeile (b) 0 $\to$ 3 von 4\,P.
(c) $x_0 = 1$ (Nullstelle der unverschobenen Logarithmusfunktion): Zeile (c) 0.}
\end{lpaufgabe}

\begin{lpaufgabe}{G8}{Logarithmen ohne Rechner}{2 P}
\begin{raster}
(a) & 1\E & $\log_2 \tfrac{1}{16} = -4$ und $\lg 0.01 = -2$. Beide nötig.\\
(b) & 1\E & $\ln e^3 = 3$ und $\log_5 1 = 0$. Beide nötig.\\
\end{raster}
\fehler{$\log_2 \tfrac{1}{16} = 4$ oder $\lg 0.01 = 2$ (Vorzeichen): Zeile (a) 0. $\log_5 1 = 1$: Zeile (b) 0. Als (E)-Weg genügt die Potenz, z.\,B. «$2^{-4} = \tfrac{1}{16}$».}
\end{lpaufgabe}

\needspace{14\baselineskip}
\section*{4 · Selbsteinschätzung}
\begin{tabularx}{\linewidth}{@{}l X@{}}\toprule
22 – 25 P & Die geprüften Teile sitzen. Wo du Punkte verloren hast: das Kapitel dieser Aufgabe nochmals (Zuordnung unten).\\
17 – 21 P & Den schwächsten Teil nochmals: Simulation und Übungen des Kapitels, in dem du die meisten Punkte verloren hast.\\
11 – 16 P & Zurück zu den Kapiteln aller Aufgaben, in denen du Punkte verloren hast.\\
0 – 10 P & Zurück zu Kapitel 1 und von dort der Reihe nach weiter.\\\midrule
\multicolumn{2}{@{}p{\linewidth}@{}}{\small\textbf{Aufgabe $\to$ Kapitel:} G1, G2 $\to$ Kapitel 1 · G3, G4 $\to$ Kapitel 2 · G5 $\to$ Kapitel 3 · G6 $\to$ Kapitel 4 · G7, G8 $\to$ Kapitel 5}\\\bottomrule
\end{tabularx}
\end{document}
